\PassOptionsToPackage{table,xcdraw}{xcolor}
\documentclass[11pt, a4paper]{article}

\usepackage[T1]{fontenc}
\usepackage[utf8]{inputenc}
\usepackage{tgpagella}
\usepackage[spanish, es-noshorthands]{babel}
\usepackage{amsmath, amssymb}
\usepackage[most]{tcolorbox}
\usepackage{tabularx}
\usepackage[margin=2.5cm]{geometry}
\usepackage{fancyhdr}

\pagestyle{fancy}
\fancyhf{}
\setlength{\headheight}{16pt}
\lhead{\textbf{UCB Sede Tarija} | Ing. Mecatrónica}
\rhead{IMT-221 | Matriz de Evidencia S08-2}
\renewcommand{\headrulewidth}{0.4pt}
\definecolor{ucbblue}{RGB}{0, 51, 102}

\begin{document}

\begin{tcolorbox}[colback=ucbblue!5, colframe=ucbblue, title=\textbf{Matriz de Evidencia Actualizada (Hito 2)}]
    Registro oficial de datos analíticos, simulados y medidos. Todos los cálculos posteriores (como $g_m$ y $A_d$) deben realizarse a partir del \textbf{Punto Q medido}, no de estimaciones nominales[cite: 5].
\end{tcolorbox}

% Se utiliza \hfill y \hrulefill para que LaTeX calcule los anchos automáticamente sin desbordar los márgenes
\noindent\textbf{Grupo:} \rule{3cm}{0.15mm} \hfill \textbf{Estado Inicio:} \rule{2.5cm}{0.15mm} \hfill \textbf{Estado Final:} \rule{2.5cm}{0.15mm} \\[0.3cm]
\noindent\textbf{Bloqueo Principal:} \hrulefill \\[0.3cm]
\noindent\textbf{Próxima Acción:} \hrulefill

\vspace{0.4cm}
\renewcommand{\arraystretch}{1.5}
\small % Se reduce la fuente levemente para acomodar las 6 columnas con holgura
\begin{tabularx}{\textwidth}{|>{\raggedright\arraybackslash}p{3.8cm}|>{\centering\arraybackslash}X|>{\centering\arraybackslash}X|>{\centering\arraybackslash}X|>{\raggedright\arraybackslash}X|>{\centering\arraybackslash}p{2.2cm}|}
\hline
\rowcolor{ucbblue!20}
\textbf{Variable} & \textbf{Predicho} & \textbf{Simulado} & \textbf{Medido} & \textbf{Interpretación} & \textbf{Status} \\ \hline
Rieles $\pm9\,\text{V}$ & & & & & \\ \hline
$I_T$ & & & & & \\ \hline
$v_E$ & & & & & \\ \hline
$v_{C1}$ & & & & & \\ \hline
$v_{C2}$ & & & & & \\ \hline
$I_{C1}$ & & & & & \\ \hline
$I_{C2}$ & & & & & \\ \hline
$V_{CE1}$ & & & & & \\ \hline
$V_{CE2}$ & & & & & \\ \hline
$V_{OD,0}$ (Offset) & & & & & \\ \hline
$g_m$ & & & & & \\ \hline
$A_d$ & & & & & \\ \hline
Resp. Diferencial ($v_{od}$) & & & & & \\ \hline
$A_c$[cite: 5] & Pendiente & & & & Pendiente \\ \hline
CMRR[cite: 5] & Pendiente & & & & Pendiente \\ \hline
\end{tabularx}
\normalsize % Restaurar el tamaño de fuente normal para el resto del documento

\vspace{0.5cm}
\begin{tcolorbox}[colback=yellow!10, colframe=orange, title=\textbf{Registro de Convención (Auditoría de Consistencia)[cite: 5]}]
    \textbf{Convención de Salida $A_d$:} \hrulefill (Ej. $v_{C2} - v_{C1}$) \\[0.2cm]
    \textbf{Convención de Salida $A_c$:} \hrulefill (Ej. $v_{C2} - v_{C1}$) \\[0.2cm]
    \textbf{¿Convención de salida consistente para CMRR? (SI / NO):} \rule{2cm}{0.15mm}
\end{tcolorbox}

\end{document}
